\section{Related Work}

We survey literature relevant to this work in three categories: RAG and its variants, attacks on RAG and Knowledge Graphs, and defenses against such attacks.

\vspace{2pt}
{\bf RAG and Variants.} 
To enhance the answer quality of LLMs on downstream tasks, methods such as agent frameworks\mcite{xia2025parallelismmeetsadaptivenessscalable,liang2024cmat,zeng2025bridgingeditinggapllms} and fine-tuning\mcite{wang2024enhancing,wang2025reasoningretrievalstudyanswer,zhu2025selfimprovingmodelsteering,yao2025countllm} are commonly used.
Among them, the RAG approach improves model responses by retrieving relevant external knowledge before generating answers\mcite{jiyuelv1,ji-etal-2024-rag,Jihealth,yi2025score,wang2025financial}.

However, conventional RAG often faces challenges such as inaccurate retrieval, hallucination during generation, and poor integration of retrieved information. 
A variety of pre-retrieval and post-retrieval strategies have been explored to address such limitations\mcite{ilin2023advanced}. Pre-retrieval improvements focus on refining indexing structures and optimizing queries through query rewriting and expansion\mcite{ma2023query, peng2024large, zheng2023take, gao2022precise}. Post-retrieval methods enhance performance through context re-ranking and compression, reducing information overload while improving coherence. 


% Further developments introduce greater adaptability and versatility by incorporating diverse strategies for improving retrieval and processing capabilities. 

%This modular architecture enables flexible component integration and reconfiguration, helping systems better handle diverse tasks and queries. Advanced techniques like RAG-Fusion and memory modules enhance retrieval by expanding query perspectives and leveraging LLM capabilities\mcite{}."

%Recent developments further improve RAG from a holistic approach. This includes adding specialized modules for search, memory, and task adaptation \mcite{yu2022generate, shao2023enhancing}. Such a modular approach allows for flexible integration and reconfiguration of components, enabling the system to better handle a variety of tasks and queries. Innovations like RAG-Fusion and the Memory module enhance retrieval by expanding query perspectives and leveraging LLM memory \mcite{wang2023knowledgpt, brown2023anchoring, cheng2024lift}. 

Recent advances take a more holistic approach to improve RAG by designing specialized modules for search, memory, and task adaptation\mcite{yu2022generate, shao2023enhancing, 
 wang2023knowledgpt, brown2023anchoring, cheng2024lift}. 
GraphRAG and its variants\mcite{graphrag, han2024retrieval,chen2025dontneedprebuiltgraphs} extends RAG by converting external knowledge into multi-scale knowledge graphs (rather than vector databases), supporting both global reasoning about broad, corpus-wide questions through community summaries and local reasoning by exploring entity relations and neighborhood structures. 


%This evolution illustrates a progression towards more sophisticated and flexible retrieval-augmented generation systems.
%Unlike RAG's reliance on semantic similarity, GraphRAG's structured knowledge representation allows for more nuanced, context-aware reasoning.


\vspace{2pt}
{\bf Attacks on RAG.}
Due to their reliance on both external knowledge bases and underlying LLMs, RAG-based models are inherently vulnerable to a variety of attacks. 
%A plethora of attacks have been proposed\mcite{zou2024poisonedrag,roychowdhury2024confusedpilot,chen2024black,tan2024glue,shafran2024machine,xue2024badrag,cheng2024trojanrag}. % These attacks involve injecting malicious data into the corpus, manipulating the retriever to return poisoned documents as relevant results for attacker-specified queries. This, in turn, increases the likelihood of LLMs generating adversarial outputs based on this poisoned information. These attacks are primarily targeted, focusing on specific data subsets and aiming to induce particular responses\mcite{zou2024poisonedrag}, rather than indiscriminately affecting the entire dataset. This targeted nature distinguishes them from untargeted attacks, which existing methods can often mitigate\mcite{zhong2023poisoning}.
% The existing attacks can be categorized according to their attack vectors:
% \ting{try to differentiate jailbreak and prompt injection attacks, seems confusing here.}
% \mct{i} Poisoning attacks inject carefully crafted poisoning texts to the knowledge base to elicit malicious responses;\mcite{zou2024poisonedrag,tan2024glue,shafran2024machine,zhong2023poisoning,}. \mcite{zou2024poisonedrag} crafts poisoning targeted answers for the specific question, thus changing the answer to the attacker-desired one.
% \mct{ii} Jailbreaking attacks aim to break the safety alignment of the LLM.  Since RAG needs to incorporate external knowledge bases, provides a special injected window for jailbreaking attacks \mcite{deng2024pandora}.
% \mct{iii} Prompt injection attacks involve directly altering the prompt given to an LLM so that the model follows the injected instructions, generating attacker-desired responses. \mcite{chen2024black} manipulate the ranking results of the retrieval model in RAG by specific instruction, \mcite{roychowdhury2024confusedpilot} embeds malicious text in the modified prompt in RAG, corrupting the responses generated by the LLM.
% \mct{iv} Backdoor attacks embed malicious functions into the LLM or knowledge base, activating such functions via specific triggers\mcite{cheng2024trojanrag, long2024backdoor,chen2024agentpoison,xue2024badrag}. \mcite{cheng2024trojanrag, xue2024badrag,chen2024agentpoison} customizes specific triggers so that when the semantics of the questions are relevant to the trigger, the answers will be impacted. \mcite{long2024backdoor} backdoor the retriever to generate the misinformation content.
% While RAG's vulnerabilities have been extensively studied, the unique security implications of GraphRAG remain largely unexplored. This work addresses this critical gap by investigating GraphRAG's unique vulnerability to poisoning attacks.
The existing attacks can be categorized according to their attack vectors. \mct{i} (Knowledge base) poisoning attacks target the knowledge base by injecting carefully crafted malicious content to manipulate RAG's responses \mcite{zou2024poisonedrag,tan2024glue,shafran2024machine,zhong2023poisoning,zhang2025benchmarkingpoisoningattacksretrievalaugmented,wang2025biasamplificationragpoisoning}. \mct{ii} Jailbreak attacks\mcite{zou2023universal,liang2025autoranweaktostrongjailbreakinglarge,jiang2024robustkv} specifically target the safety guardrails of RAG's underlying LLMs. Notably, while typical jailbreak attacks target LLM safety guardrails directly, RAG models are particularly vulnerable because their external knowledge bases create additional attack surfaces\mcite{deng2024pandora,wang2025privacyawaredecodingmitigatingprivacy}. \mct{iii} Prompt injection attacks manipulate input prompts to override intended system behavior. In the RAG context, these attacks operate through two mechanisms: manipulating retrieval rankings via specific instructions\mcite{chen2024black,zuo2025makemedicalaisystems,dai2025disablingselfcorrectionretrievalaugmentedgeneration}, and embedding malicious content within modified prompts to corrupt generated responses\mcite{roychowdhury2024confusedpilot}. \mct{iv} Backdoor attacks embed malicious functionalities into RAG models that are activated through specific triggers\mcite{cheng2024trojanrag, long2024backdoor,chen2024agentpoison,xue2024badrag,bagwe2025ragunfairexposingfairness}. Instances of such attacks include semantic triggers that respond to specific question content\mcite{cheng2024trojanrag, xue2024badrag,chen2024agentpoison} or retriever-level backdoors that generate targeted misinformation\mcite{long2024backdoor}.


While these vulnerabilities have been extensively studied in the context of conventional RAG, the security implications of GraphRAG remain largely unexplored. This work bridges this critical gap by examining GraphRAG's unique vulnerabilities to knowledge poisoning attacks.

%All these attacks need to inject poisoned data into the underlying data corpus. This makes them essentially targeted data poisoning attacks against the retrievers.
\vspace{2pt}
{\bf Attacks on Knowledge Graphs.} Zhang et al.\mcite{zhangDataPoisoningAttack2019} highlight the susceptibility of knowledge graph embedding models to data poisoning, demonstrating that manipulating a small number of triples can significantly alter link prediction. Subsequent work explores more targeted poisoning strategies\mcite{bhardwajPoisoningKnowledgeGraph2021} and reveals the vulnerabilities of KG-based recommender systems\mcite{wuPoisoningAttacksKnowledge2022} and federated learning\mcite{zhouPoisoningAttackFederated2024}. More recently, Xi et al.\mcite{xi2023security} introduce a poisoning attack designed to hijack KG-based reasoning queries without impacting non-target performance. However, these attacks are inapplicable for \rag with LLMs, due to their reliance on text embeddings and specific retrieval mechanisms.

\vspace{2pt}
{\bf Defenses against RAG Poisoning Attacks.}
Prior work has proposed defenses such as perplexity-based detection~\mcite{jelinek1980interpolated}, query paraphrasing~\mcite{jain2023baseline}, and expanded context windows~\mcite{shafran2024machine}. However, these techniques have shown limited effectiveness~\mcite{chen2024agentpoison,zou2024poisonedrag}, as they do not address the core vulnerability: the retrieval corpus's susceptibility to targeted poisoning.

Recent work explores more advanced defenses that target RAG's fundamental vulnerability by filtering malicious content and reconciling conflicts between the LLM's internal knowledge and retrieved external information. For instance, TrustRAG\mcite{zhou2025trustrag} employs K-means clustering to filter malicious data and resolve knowledge conflicts. RobustRAG\mcite{xiang2024certifiably} introduces an `isolate-then-aggregate'' framework that generates responses from individual passages before secure aggregation, providing certifiable robustness for certain queries. AstuteRAG~\mcite{wang2024astuteragovercomingimperfect} iteratively combines internal and external knowledge with source-aware filtering.

% Recent advances aim to mitigate this root cause by filtering malicious content and reconciling conflicts between the LLM’s internal knowledge and retrieved context. For example, TrustRAG~\mcite{zhou2025trustrag} uses K-means clustering for conflict resolution, RobustRAG~\mcite{xiang2024certifiably} introduces an ``isolate-then-aggregate'' mechanism for certifiable robustness, and AstuteRAG~\mcite{wang2024astuteragovercomingimperfect} iteratively combines internal and external knowledge with source-aware filtering.

However, these defenses are not directly applicable to GraphRAG due to its complex context construction, where multiple entities, relations, summaries, and text chunks are interwoven, rendering context segmentation and passage-level filtering ineffective. Our work thus explores defenses specifically tailored to GraphRAG and their inherent limitations.



\section{Conclusion}



We present a systematic study of GraphRAG's unique vulnerabilities to poisoning attacks. Our analysis reveals a critical security paradox: \jc{while the graph-based indexing and retrieval pipeline in GraphRAG reduces the effectiveness of existing RAG poisoning attacks,} these features also introduce new attack surfaces. Specifically, the adversary can exploit the knowledge graph structure to craft poisoning text targeting multiple queries simultaneously, enabling more effective and scalable attacks. We further examine the unique challenges of defending against such attacks, identifying several promising directions for future research.


\section*{Acknowledgements}
\jcc{This work is supported by the National Science Foundation under Grant No. 2405136 and 2406572, and OpenAI's Researcher Access Program. }